% =====================================================================
%  Лекция 6. Точечные оценки. Свойства оценок
%  Дата: 2026-10-22
% =====================================================================
\section{Точечные оценки. Свойства оценок}
\label{lec:06}

\subsection{Постановка задачи}
\label{lec06:setting}

Пусть задана случайная величина~$\xi$, функция распределения которой известна с точностью
до параметра:
\[
  F_\xi(x) = F_\xi(x, \theta), \qquad \theta = (\theta_1, \ldots, \theta_m) \in \Th \subset \R^m,
\]
где $\Th$~--- \emph{параметрическое множество}, множество всех допустимых значений параметра.
Сам параметр неизвестен, и его нужно оценить по выборке $X_{[n]} = (X_1, \ldots, X_n)$
из генеральной совокупности~$\xi$.

\begin{example}
  Нормальное распределение $\Norm(a, \sigma^2)$:
  \[
    F_\xi(x, \theta) = \int_{-\infty}^{x} \frac{1}{\sqrt{2\pi}\,\sigma}
                       e^{-\frac{(y - a)^2}{2\sigma^2}} \, dy, \qquad
    \theta = (a, \sigma^2) \in \Th = \R \times (0, +\infty).
  \]
  Здесь два неизвестных параметра, $m = 2$.
\end{example}

\begin{definition}
  \emph{Оценкой} неизвестного параметра~$\theta$ называется статистика
  $\hat\theta(X_{[n]})$ (то есть борелевская функция от выборки), значение которой
  рассматривается как приближённое значение~$\theta$.
\end{definition}

Под это определение подходит любая статистика, так что оценок одного параметра можно
придумать сколько угодно. Нужны критерии, по которым оценки сравнивают,~--- \emph{свойства
оценок}.

\subsection{Несмещённость}
\label{lec06:unbiased}

\begin{definition}
  Оценка $\hat\theta(X_{[n]})$ называется \emph{несмещённой} оценкой параметра~$\theta$, если
  \[
    \E \hat\theta(X_{[n]}) = \theta \qquad \forall\, \theta \in \Th.
  \]
\end{definition}

Смысл: оценка не ошибается <<систематически>>~--- в среднем она даёт верное значение,
каким бы ни было истинное~$\theta$.

\begin{definition}
  Оценка $\hat\theta(X_{[n]})$ называется \emph{асимптотически несмещённой} оценкой
  параметра~$\theta$, если
  \[
    \E \hat\theta(X_{[n]}) \xrightarrow[n \to \infty]{} \theta \qquad \forall\, \theta \in \Th.
  \]
\end{definition}

\begin{example}[выборочное среднее]
  \label{lec06:ex:mean}
  Пусть $\E\xi$ существует и конечно. Так как $X_i$ распределены как~$\xi$,
  \[
    \E\Xbar = \E\biggl(\frac1n \sum_{i=1}^n X_i\biggr) = \frac1n \sum_{i=1}^n \E X_i = \E\xi.
  \]
  Значит, $\Xbar$~--- несмещённая оценка параметра $\theta = \E\xi$.
\end{example}

\begin{example}[выборочная дисперсия]
  \label{lec06:ex:var}
  Пусть $\theta = \D\xi < \infty$. Проверим, будет ли несмещённой
  $s^2 = \frac1n \sum_{i=1}^n (X_i - \Xbar)^2$.

  Прибавим и вычтем $\E\xi$ под квадратом и раскроем его:
  \begin{align*}
    s^2 &= \frac1n \sum_{i=1}^n \bigl((X_i - \E\xi) - (\Xbar - \E\xi)\bigr)^2 \\
        &= \frac1n \sum_{i=1}^n (X_i - \E\xi)^2
           - \frac2n (\Xbar - \E\xi) \sum_{i=1}^n (X_i - \E\xi)
           + (\Xbar - \E\xi)^2.
  \end{align*}
  Множитель $(\Xbar - \E\xi)$ от $i$ не зависит, поэтому вынесен за сумму. Далее
  $\sum_{i=1}^n (X_i - \E\xi) = n\Xbar - n\E\xi = n(\Xbar - \E\xi)$, так что средний член
  равен $-2(\Xbar - \E\xi)^2$ и
  \[
    s^2 = \frac1n \sum_{i=1}^n (X_i - \E\xi)^2 - (\Xbar - \E\xi)^2.
  \]
  Берём математическое ожидание. Так как $\E X_i = \E\xi$ и $\E\Xbar = \E\xi$
  (пример~\ref{lec06:ex:mean}), это дисперсии:
  \[
    \E s^2 = \frac1n \sum_{i=1}^n \D X_i - \D\Xbar = \D\xi - \D\Xbar.
  \]
  Осталось найти $\D\Xbar$. Дисперсия суммы \emph{независимых} величин равна сумме
  дисперсий (для зависимых добавились бы ковариации), поэтому
  \[
    \D\Xbar = \D\biggl(\frac1n \sum_{i=1}^n X_i\biggr) = \frac{1}{n^2} \sum_{i=1}^n \D X_i
            = \frac{1}{n^2} \cdot n\,\D\xi = \frac{\D\xi}{n}.
  \]
  Итого
  \[
    \E s^2 = \D\xi - \frac1n \D\xi = \frac{n-1}{n}\, \D\xi \ne \D\xi.
  \]
  Значит, $s^2$~--- \emph{смещённая} оценка дисперсии (в среднем она её занижает).
  Но $\frac{n-1}{n} \to 1$, поэтому $s^2$~--- асимптотически несмещённая оценка.
\end{example}

Смещение легко убрать, домножив на $\frac{n}{n-1}$:

\begin{definition}
  \emph{Исправленной выборочной дисперсией} называется
  \[
    \tilde s^2 = \frac{n}{n-1}\, s^2 = \frac{1}{n-1} \sum_{i=1}^n (X_i - \Xbar)^2.
  \]
\end{definition}
По линейности $\E\tilde s^2 = \frac{n}{n-1} \cdot \frac{n-1}{n}\, \D\xi = \D\xi$, то есть
$\tilde s^2$~--- несмещённая оценка дисперсии.

\subsection{Состоятельность}
\label{lec06:consistency}

Далее $\theta \in \Th \subset \R$.

\begin{reminder}[Сходимость по вероятности и почти наверное]
  \begin{itemize}
    \item $\xi_n \toP \xi$, если $\P\{\abs{\xi_n - \xi} \ge \epsilon\} \to 0$
          при $n \to \infty$ для любого $\epsilon > 0$;
    \item $\xi_n \toas \xi$, если $\P\bigl\{\omega : \xi_n(\omega) \to \xi(\omega)\bigr\} = 1$.
  \end{itemize}
  Из сходимости п.н.\ следует сходимость по вероятности, а из неё~--- сходимость
  по распределению.
\end{reminder}

\begin{definition}
  Оценка $\hat\theta(X_{[n]})$ называется
  \begin{itemize}
    \item \emph{состоятельной} оценкой параметра~$\theta$, если
          $\hat\theta(X_{[n]}) \toP \theta$ при $n \to \infty$ для любого $\theta \in \Th$;
    \item \emph{сильно состоятельной}, если
          $\hat\theta(X_{[n]}) \toas \theta$ при $n \to \infty$ для любого $\theta \in \Th$.
  \end{itemize}
\end{definition}

Сильно состоятельная оценка состоятельна (п.н.\ $\Rightarrow$ по вероятности).

Смысл: с ростом объёма выборки оценка приближается к истинному значению параметра.

\subsubsection{Состоятельность и статистики первого типа}

Сильную состоятельность многих оценок мы фактически уже доказали. По
теореме~\ref{lec03:first-type-conv} статистика первого типа сходится п.н.:
\[
  S(X_{[n]}) = G(F_n^*) \toas G(F_\xi) = h(a).
\]
Поэтому если оценка~--- статистика первого типа и $h(a)$ совпадает с оцениваемым параметром,
то она сильно состоятельна. Например, при условии, что соответствующие моменты существуют
и конечны:
\[
  \Xbar = a_1^* \toas a_1 = \E\xi, \qquad
  a_k^* \toas a_k, \qquad
  a_k^{0*} \toas a_k^0, \qquad
  s^2 = a_2^{0*} \toas a_2^0 = \D\xi.
\]
Так же сходятся к своим теоретическим аналогам выборочные коэффициенты асимметрии
и эксцесса (раздел~\ref{lec03:moments}).

\begin{remark}
  Выборочная дисперсия~--- это \emph{центральный} момент $a_2^{0*}$. Начальный момент
  $a_2^* = \frac1n \sum X_i^2$ сходится к $\E\xi^2$, а не к дисперсии.
\end{remark}

\subsection{Эффективность}
\label{lec06:efficiency}

Пусть $\theta \in \Th \subset \R$ и задан некоторый класс оценок~$K$ параметра~$\theta$.

\begin{definition}
  Оценка $\theta^*(X_{[n]}) \in K$ называется \emph{эффективной} оценкой параметра~$\theta$
  в классе~$K$, если для любой другой оценки $\hat\theta \in K$
  \begin{equation}
    \label{lec06:eq:eff}
    \E(\theta^* - \theta)^2 \le \E(\hat\theta - \theta)^2 \qquad \forall\, \theta \in \Th.
  \end{equation}
\end{definition}

$\E(\hat\theta - \theta)^2$~--- среднеквадратичная ошибка оценки. Эффективная оценка~---
та, у которой она наименьшая в классе~$K$, причём при \emph{каждом} значении~$\theta$.

Важный частный случай: $K$~--- класс \emph{несмещённых} оценок, $\E\hat\theta = \theta$
для всех $\hat\theta \in K$. Тогда $\theta = \E\theta^* = \E\hat\theta$, и
неравенство~\eqref{lec06:eq:eff} принимает вид
\[
  \E(\theta^* - \E\theta^*)^2 \le \E(\hat\theta - \E\hat\theta)^2,
  \qquad\text{то есть}\qquad
  \D\theta^* \le \D\hat\theta.
\]
Эффективная несмещённая оценка~--- это несмещённая оценка с наименьшей дисперсией.

\begin{theorem}[{\cite[гл.~15, теорема~1.1]{bure}}]
  \label{lec06:uniqueness}
  Пусть несмещённые оценки $\hat\theta_1$ и $\hat\theta_2$ параметра $\theta \in \Th \subset \R$
  являются эффективными (в классе несмещённых оценок). Тогда $\hat\theta_1$ и $\hat\theta_2$
  почти наверное совпадают.
\end{theorem}

\begin{proof}
  Обе оценки несмещённые и эффективные, поэтому
  \[
    \E\hat\theta_1 = \E\hat\theta_2 = \theta, \qquad \D\hat\theta_1 = \D\hat\theta_2
  \]
  (каждая из дисперсий не больше другой). Рассмотрим полусумму
  \[
    \tilde\theta = \frac{\hat\theta_1 + \hat\theta_2}{2}, \qquad
    \E\tilde\theta = \frac{\theta + \theta}{2} = \theta,
  \]
  она тоже несмещённая. Справедливо алгебраическое тождество (проверяется раскрытием скобок,
  это $\frac{u^2 + v^2}{2} = \bigl(\frac{u - v}{2}\bigr)^2 + \bigl(\frac{u + v}{2}\bigr)^2$
  при $u = \hat\theta_1 - \theta$, $v = \hat\theta_2 - \theta$):
  \[
    \frac{(\hat\theta_1 - \theta)^2}{2} + \frac{(\hat\theta_2 - \theta)^2}{2}
    = \biggl(\frac{\hat\theta_1 - \hat\theta_2}{2}\biggr)^{\!2} + (\tilde\theta - \theta)^2.
  \]
  Возьмём математическое ожидание. Для несмещённых оценок $\E(\hat\theta - \theta)^2 = \D\hat\theta$,
  поэтому
  \[
    \D\hat\theta_1 = \frac{\D\hat\theta_1}{2} + \frac{\D\hat\theta_2}{2}
    = \frac14\, \E(\hat\theta_1 - \hat\theta_2)^2 + \D\tilde\theta.
  \]
  Оценка $\tilde\theta$ несмещённая, а $\hat\theta_1$ эффективна, значит,
  $\D\tilde\theta \ge \D\hat\theta_1$. Отсюда
  \[
    \frac14\, \E(\hat\theta_1 - \hat\theta_2)^2 = \D\hat\theta_1 - \D\tilde\theta \le 0,
  \]
  а так как левая часть неотрицательна, $\E(\hat\theta_1 - \hat\theta_2)^2 = 0$.
  Неотрицательная величина с нулевым средним равна нулю п.н., поэтому
  $\hat\theta_1 = \hat\theta_2$ п.н.
\end{proof}

Иначе говоря, эффективная несмещённая оценка, если существует, единственна.

\subsection{Асимптотическая нормальность}
\label{lec06:asymptotic-normality}

Пусть $\theta \in \Th \subset \R$.

\begin{definition}
  Оценка $\hat\theta(X_{[n]})$ называется \emph{асимптотически нормальной} оценкой
  параметра~$\theta$ с коэффициентом (рассеивания) $\sigma^2(\theta)$, если
  \[
    \sqrt{n}\,\bigl(\hat\theta(X_{[n]}) - \theta\bigr) \tod \zeta \sim \Norm\bigl(0, \sigma^2(\theta)\bigr)
    \qquad \text{при } n \to \infty, \ \forall\, \theta \in \Th.
  \]
\end{definition}

\begin{reminder}[Сходимость по распределению]
  $\xi_n \tod \xi$, если $F_{\xi_n}(x) \to F_\xi(x)$ в каждой точке~$x$, где $F_\xi$ непрерывна
  (лекция~\ref{lec:02}). Равносильное определение: $\E g(\xi_n) \to \E g(\xi)$ для любой
  ограниченной непрерывной функции~$g$.
\end{reminder}

У нормального закона функция распределения непрерывна всюду, поэтому определение
означает: для любого $x \in \R$
\[
  \P\bigl\{\sqrt{n}\,(\hat\theta - \theta) \le x\bigr\} \xrightarrow[n \to \infty]{}
  \frac{1}{\sqrt{2\pi}\,\sigma(\theta)} \int_{-\infty}^{x} e^{-\frac{y^2}{2\sigma^2(\theta)}} \, dy.
\]
Если $\zeta \sim \Norm(a, \sigma^2)$, то $\frac{\zeta - a}{\sigma} \sim \Norm(0, 1)$. Поэтому
определение можно записать и через стандартный нормальный закон:
\[
  \frac{\sqrt{n}\,(\hat\theta - \theta)}{\sigma(\theta)} \tod \zeta' \sim \Norm(0, 1).
\]

\subsubsection{Сравнение асимптотически нормальных оценок}

Пусть $\hat\theta_1$ и $\hat\theta_2$~--- асимптотически нормальные оценки~$\theta$
с коэффициентами $\sigma_1^2(\theta)$ и $\sigma_2^2(\theta)$, и $\sigma_1^2(\theta) < \sigma_2^2(\theta)$.
Нормированные величины
\[
  \frac{\sqrt{n}\,(\hat\theta_1 - \theta)}{\sigma_1(\theta)}
  \qquad\text{и}\qquad
  \frac{\sqrt{n}\,(\hat\theta_2 - \theta)}{\sigma_2(\theta)}
\]
имеют в пределе одно и то же распределение $\Norm(0, 1)$. Значит, при больших~$n$
приближённо $\hat\theta_k \approx \theta + \frac{\sigma_k(\theta)}{\sqrt{n}}\, \zeta'$:
обе оценки разбросаны вокруг~$\theta$ по нормальному закону, но у $\hat\theta_1$ разброс меньше.
Поэтому $\hat\theta_1$ лучше: из двух асимптотически нормальных оценок предпочтительнее
та, у которой меньше коэффициент $\sigma^2(\theta)$.

\begin{keybox}[Что нужно знать к экзамену]
  \begin{itemize}
    \item Определение оценки; параметрическое множество~$\Th$.
    \item Несмещённость и асимптотическая несмещённость; $\E\Xbar = \E\xi$.
    \item Вывод $\E s^2 = \frac{n-1}{n} \D\xi$ (через $\D\Xbar = \D\xi / n$);
          исправленная дисперсия $\tilde s^2$ и её несмещённость.
    \item Состоятельность и сильная состоятельность; почему статистики первого типа
          дают сильно состоятельные оценки (моменты, $\Xbar$, $s^2$).
    \item Эффективность в классе~$K$; для несмещённых оценок~--- минимальная дисперсия.
    \item Теорема о единственности эффективной несмещённой оценки с доказательством.
    \item Асимптотическая нормальность, коэффициент $\sigma^2(\theta)$; как сравнивать
          асимптотически нормальные оценки.
  \end{itemize}
  \emph{Где легко ошибиться:}
  \begin{itemize}
    \item $s^2$ смещена, $\tilde s^2$ несмещена; в $\tilde s^2$ делим на $n - 1$;
    \item $s^2 = a_2^{0*}$~--- центральный момент; $a_2^* \to \E\xi^2$, а не $\D\xi$;
    \item $\D\Xbar = \D\xi / n$ использует независимость $X_i$;
    \item сходимость п.н.: вероятность множества, где $\xi_n \to \xi$, равна~$1$, а не~$0$;
    \item в доказательстве единственности нужен именно $\E(\hat\theta_1 - \hat\theta_2)^2 = 0$
          (с квадратом), отсюда равенство п.н.;
    \item в эффективности неравенство должно выполняться при всех $\theta \in \Th$.
  \end{itemize}
\end{keybox}
